\documentclass[twoside,12pt]{book}

\usepackage[utf8]{inputenc}
\usepackage[T1]{fontenc}
\usepackage{mathpazo}
\usepackage{amsmath,amssymb,amsthm}
\usepackage{graphicx}
\usepackage{listings}
\usepackage{xcolor}
\usepackage{enumitem}
\usepackage{tikz}
\usepackage{algorithm}
\usepackage{algpseudocode}
\usepackage{booktabs}
\usepackage{geometry}
\usepackage{fancyhdr}
\usepackage{titlesec}
\usepackage[colorlinks=true, linkcolor=blue, citecolor=blue, urlcolor=blue, plainpages=false, pdfpagelabels]{hyperref}

\geometry{
  papersize={6in,9in},
  margin=0.75in,
  top=0.75in,bottom=0.75in,
  inner=0.85in,outer=0.65in,
  headheight=14pt
}

\pagestyle{fancy}
\fancyhf{}
\fancyhead[LE]{\small\textit{6. The Probabilistic Method}}
\fancyhead[RO]{\small\textit{Randomized Algorithms}}
\fancyfoot[C]{\thepage}
\renewcommand{\headrulewidth}{0.4pt}

\definecolor{codegray}{rgb}{0.45,0.45,0.45}
\lstdefinestyle{cppstyle}{
  basicstyle=\ttfamily\scriptsize,
  keywordstyle=\bfseries,
  commentstyle=\itshape\color{codegray},
  numbers=left,
  numberstyle=\tiny\color{codegray},
  numbersep=8pt,
  frame=single,
  framerule=0.4pt,
  rulecolor=\color{codegray},
  backgroundcolor=\color{gray!5},
  breaklines=true,
  tabsize=4,
  showstringspaces=false,
  language=C++,
  morekeywords={nullptr,size_t,auto,const,constexpr,
    unordered_map,vector,pair,mt19937,
    uniform_int_distribution,INT_MAX}
}
\lstset{style=cppstyle}

\theoremstyle{plain}
\newtheorem{theorem}{Theorem}[chapter]
\newtheorem{lemma}[theorem]{Lemma}
\newtheorem{corollary}[theorem]{Corollary}
\newtheorem{proposition}[theorem]{Proposition}
\theoremstyle{definition}
\newtheorem{definition}[theorem]{Definition}
\newtheorem{example}[theorem]{Example}
\newtheorem{remark}[theorem]{Remark}
\theoremstyle{remark}
\newtheorem{exercise}[theorem]{Exercise}
\newtheorem{problem}[theorem]{Problem}
\newenvironment{cpplisting}{\begin{center}\small}{\end{center}}

\newcommand{\Exp}{\operatorname{\mathbb{E}}}
\newcommand{\Var}{\operatorname{Var}}
\newcommand{\Cov}{\operatorname{Cov}}

\begin{document}

\chapter*{6 The Probabilistic Method}
\addcontentsline{toc}{chapter}{6 The Probabilistic Method}
\markboth{6 The Probabilistic Method}{Randomized Algorithms}
\label{ch:probabilistic-method}

The probabilistic method is one of the most powerful and elegant techniques
in combinatorics and theoretical computer science. Introduced by
Erd\H{o}s in the 1940s, its core idea is deceptively simple: to prove that
a combinatorial object with certain properties \emph{exists}, it suffices
to show that a randomly constructed object has the desired properties with
positive probability. Unlike constructive proofs, the probabilistic method
need not exhibit the object explicitly---it merely demonstrates that the
object cannot fail to exist.

The method has three essential ingredients. First, define a probability
space over candidate objects. Second, identify the ``bad'' events---those
properties that the object must avoid. Third, show that the probability of
any bad event is strictly less than one, so that with positive probability
no bad event occurs. The power of the method lies in the fact that the
probabilistic analysis often reveals structural truths that are invisible
to purely algebraic or constructive approaches.

In this chapter, we develop the probabilistic method from its foundations
through several important applications: random satisfiability, expander
graphs, oblivious routing, the Lov\'{a}sz Local Lemma, and the method of
conditional probabilities for derandomization.

%======================================================================
\section{Overview of the Method}
\label{sec:prob-method-overview}
%======================================================================

The simplest form of the probabilistic method is captured by the following
observation. Suppose we wish to show that there exists an object $x$ in a
finite set $S$ satisfying property $P$. If we define a probability
distribution on $S$ and show that
\[
\Pr_{x \leftarrow S}[P(x)] > 0,
\]
then there must exist at least one $x \in S$ with property $P$.

More generally, suppose we have $m$ ``bad'' events $A_1, A_2, \ldots, A_m$
that we wish to avoid simultaneously. If we can show that
\[
\Pr\!\Bigl[\bigcup_{i=1}^{m} A_i\Bigr] < 1,
\]
then $\Pr[\overline{A_1} \cap \overline{A_2} \cap \cdots \cap \overline{A_m}]
> 0$, and the desired object exists. The union bound (Boole's inequality)
gives the simple sufficient condition
\[
\sum_{i=1}^{m} \Pr[A_i] < 1.
\]
More refined versions, such as the Lov\'{a}sz Local Lemma (Section~\ref{sec:lll}),
provide much sharper conditions.

\begin{example}[Ramsey Numbers]
\label{ex:ramsey}
The Ramsey number $R(3,3)$ is the minimum $n$ such that every 2-coloring
of the edges of the complete graph $K_n$ contains a monochromatic triangle.
We prove the upper bound $R(3,3) \leq 6$ using the probabilistic method
(though a direct argument suffices for this small case).

More interestingly, consider the question: for a given $n$, what is the
probability that a random 2-coloring of $K_n$ avoids monochromatic
triangles? Color each edge independently red or blue with probability
$1/2$ each. For a fixed triangle, the probability that all three edges
receive the same color is $2 \cdot (1/2)^3 = 1/4$. By linearity of
expectation, the expected number of monochromatic triangles is
\[
\binom{n}{3} \cdot \frac{1}{4}.
\]
When $n = 5$, this expectation is $\binom{5}{3}/4 = 10/4 = 2.5$. When
$n = 4$, it is $\binom{4}{3}/4 = 4/4 = 1$. The probabilistic argument
shows that for $n = 5$, since the expected count is $2.5 > 0$, there
exist colorings with monochromatic triangles---but this does not yet show
that \emph{every} coloring has one. The power of the method is more
evident in larger Ramsey problems where direct enumeration is infeasible.
\end{example}

\begin{lemma}[Expected Monochromatic Subgraphs]
\label{lem:monochromatic}
For a random 2-coloring of the edges of $K_n$ (each edge colored independently
red or blue with probability $1/2$), the expected number of monochromatic
copies of $K_k$ is
\[
\binom{n}{k} \cdot 2^{1 - \binom{k}{2}}.
\]
\end{lemma}

\begin{proof}
There are $\binom{n}{k}$ copies of $K_k$ in $K_n$. For a fixed copy, the
probability that all $\binom{k}{2}$ edges receive the same color is
$2 \cdot (1/2)^{\binom{k}{2}} = 2^{1 - \binom{k}{2}}$. By linearity of
expectation, the expected number of monochromatic copies is the product.
\end{proof}

This simple counting argument---relying only on linearity of
expectation---yields non-trivial upper bounds on Ramsey numbers and is
the prototype for the probabilistic method.

%======================================================================
\section{Maximum Satisfiability}
\label{sec:maxsat}
%======================================================================

The Maximum Satisfiability problem (MAX-SAT) asks: given a Boolean
formula in conjunctive normal form (CNF), find an assignment that
satisfies the maximum number of clauses. This problem is NP-hard in
general, but the probabilistic method provides elegant approximation
guarantees.

\begin{definition}[MAX-SAT]
Let $\phi = C_1 \wedge C_2 \wedge \cdots \wedge C_m$ be a CNF formula
with $n$ variables $x_1, \ldots, x_n$ and $m$ clauses. Each clause $C_j$
is a disjunction of literals. The MAX-SAT problem asks for an assignment
$\alpha : \{x_1, \ldots, x_n\} \to \{0, 1\}$ maximizing the number of
satisfied clauses $|\{j : \alpha \models C_j\}|$.
\end{definition}

\begin{theorem}[Random Assignment for MAX-SAT]
\label{thm:maxsat-random}
For any CNF formula $\phi$ with $m$ clauses, a random assignment satisfies
at least $m/2$ clauses in expectation. Consequently, there exists an
assignment satisfying at least $m/2$ clauses.
\end{theorem}

\begin{proof}
Let $X_j$ be the indicator random variable for the event that clause $C_j$
is satisfied by a random assignment (each variable set to $0$ or $1$
independently with probability $1/2$). Since $C_j$ is a disjunction of
$l_j$ literals, the probability that $C_j$ is \emph{not} satisfied is
$(1/2)^{l_j}$. Thus
\[
\mathbb{E}[X_j] = \Pr[C_j \text{ is satisfied}] = 1 - 2^{-l_j} \geq 1 - \frac{1}{2} = \frac{1}{2}
\]
for every clause with at least one literal. By linearity of expectation,
\[
\mathbb{E}\!\Bigl[\sum_{j=1}^{m} X_j\Bigr] = \sum_{j=1}^{m} \mathbb{E}[X_j]
\geq \frac{m}{2}.
\]
Since the expected number of satisfied clauses is at least $m/2$, there
must exist an assignment satisfying at least $m/2$ clauses.
\end{proof}

This result extends immediately to MAX-3SAT, where each clause has exactly
three literals.

\begin{corollary}[7/8-Approximation for MAX-3SAT]
\label{cor:max3sat}
For any 3-CNF formula $\phi$ with $m$ clauses, a random assignment satisfies
at least $7m/8$ clauses in expectation. Hence MAX-3SAT has a
$7/8$-approximation.
\end{corollary}

\begin{proof}
Each clause has exactly three literals, so
\[
\mathbb{E}[X_j] = 1 - 2^{-3} = 1 - \frac{1}{8} = \frac{7}{8}.
\]
By linearity of expectation, $\mathbb{E}[\sum X_j] = 7m/8$.
\end{proof}

The ratio $7/8$ is the best achievable by any random assignment analysis
for MAX-3SAT. Tighter bounds require more sophisticated techniques, but
this simple argument demonstrates the power of the probabilistic method
for approximation algorithms.

%======================================================================
\section{Expanding Graphs}
\label{sec:expanders}
%======================================================================

Expander graphs are sparse graphs with strong connectivity properties:
every small set of vertices has many neighbors. They arise naturally in
the probabilistic method as objects that exist with high probability.

\begin{definition}[Expander Graph]
A $d$-regular graph $G = (V, E)$ on $n$ vertices is an \emph{expander} if
its second-largest eigenvalue $\lambda_2$ of the adjacency matrix satisfies
\[
\lambda_2 \leq 2\sqrt{d - 1}.
\]
Such graphs are called \emph{Ramanujan graphs} when this bound is achieved.
\end{definition}

The spectral gap $\lambda_1 - \lambda_2 = d - \lambda_2$ controls the
expansion quality. By the Cheeger inequality, the edge expansion $h(G)$
is related to the spectral gap by
\[
\frac{d - \lambda_2}{2} \leq h(G) \leq \sqrt{2d(d - \lambda_2)}.
\]

\begin{theorem}[Random Regular Graphs are Expanders]
\label{thm:random-expander}
For every fixed degree $d \geq 3$, a random $d$-regular graph on $n$
vertices has $\lambda_2 \leq 2\sqrt{d-1} + \epsilon$ with high probability
for any $\epsilon > 0$, as $n \to \infty$.
\end{theorem}

The proof uses the trace method: the $k$-th moment of
$\operatorname{tr}(A^k)$ counts closed walks of length $k$ in $G$, and
concentration of this count yields concentration of the eigenvalue
distribution. For Ramanujan-quality expanders, explicit constructions
based on number theory (Lubotzky--Phillips--Sarnak, Margulis) achieve the
bound $\lambda_2 \leq 2\sqrt{d-1}$ deterministically.

The following lemma characterizes the spectral gap in terms of graph
connectivity.

\begin{lemma}[Cheeger Inequality]
\label{lem:cheeger}
For a $d$-regular graph $G$ on $n$ vertices with edge expansion
\[
h(G) = \min_{S \subseteq V, |S| \leq n/2} \frac{|\partial S|}{|S|},
\]
we have
\[
\frac{d - \lambda_2}{2} \leq h(G) \leq \sqrt{2d(d - \lambda_2)}.
\]
\end{lemma}

%======================================================================
\section{Oblivious Routing Revisited}
\label{sec:oblivious-routing}
%======================================================================

In Chapter~5, we analyzed random routing on the complete graph. We now
consider oblivious routing on general graphs, where the routing strategy
must be fixed in advance without knowledge of the traffic pattern.

\begin{definition}[Oblivious Routing]
An oblivious routing strategy on a graph $G$ is a probability distribution
over paths for each source-destination pair $(s, t)$, specified without
knowledge of the other traffic demands. The \emph{competitive ratio} is
the worst-case ratio of the expected maximum congestion to the optimal
offline maximum congestion.
\end{definition}

\begin{theorem}[Random Oblivious Routing]
\label{thm:oblivious-routing}
On any graph $G$ with $n$ vertices and average degree $\bar{d}$, the
random oblivious routing strategy---choosing a shortest path uniformly at
random for each demand---achieves competitive ratio $O(\log n)$ with high
probability.
\end{theorem}

The analysis proceeds by decomposing each shortest path into edges, then
applying the Chernoff bound to the congestion at each edge. For each edge
$e$, the congestion $X_e$ is the sum of independent indicator variables
(one per path using $e$), and standard concentration inequalities yield
the $O(\log n)$ bound on the maximum congestion.

The key insight connecting this to the probabilistic method is that the
\emph{existence} of a good routing is guaranteed by the probabilistic
argument: the expected maximum congestion is $O(\log n)$, so a routing
with maximum congestion $O(\log n)$ exists.

%======================================================================
\section{The Lov\'{a}sz Local Lemma}
\label{sec:lll}
%======================================================================

The union bound $\Pr[\bigcup A_i] \leq \sum \Pr[A_i]$ is often too crude
when the bad events are not ``spread out'' across the probability space.
The Lov\'{a}sz Local Lemma provides a much sharper tool when the bad
events have limited dependencies.

\begin{definition}[Dependency Graph]
Let $A_1, \ldots, A_m$ be events in a probability space. A
\emph{dependency graph} for these events is a graph $H$ on vertex set
$\{1, \ldots, m\}$ such that for each $i$, the event $A_i$ is independent
of the events $\{A_j : j \notin N_H(i)\}$, where $N_H(i)$ denotes the
neighbors of $i$ in $H$.
\end{definition}

\begin{theorem}[Lov\'{a}sz Local Lemma---Asymmetric Form]
\label{thm:lll}
Let $A_1, A_2, \ldots, A_m$ be events with a dependency graph $H$ of
maximum degree $d$. If there exist real numbers $x_1, x_2, \ldots, x_m$
in $[0, 1)$ such that
\[
\Pr[A_i] \leq x_i \prod_{j \in N_H(i)} (1 - x_j)
\qquad \text{for all } i = 1, \ldots, m,
\]
then
\[
\Pr\!\Bigl[\bigcap_{i=1}^{m} \overline{A_i}\Bigr]
> \prod_{i=1}^{m} (1 - x_i) > 0.
\]
\end{theorem}

The symmetric form is often more convenient.

\begin{corollary}[Symmetric Local Lemma]
\label{cor:lll-symmetric}
If each event $A_i$ has $\Pr[A_i] \leq p$ and depends on at most $d$
other events, and if
\[
ep(d + 1) \leq 1,
\]
then $\Pr[\bigcap \overline{A_i}] > 0$.
\end{corollary}

\begin{proof}
Set $x_i = 1/(d+1)$ for all $i$. The condition $\Pr[A_i] \leq
x_i(1 - x_i)^d$ is implied by $p \leq \frac{1}{e(d+1)}$ since
$(1 - 1/(d+1))^d \geq 1/e$.
\end{proof}

\begin{example}[2-SAT]
Consider a 2-SAT formula $\phi$ on $n$ variables with $m$ clauses. Each
clause has two literals. Assign each variable independently to $0$ or $1$
with probability $1/2$. The bad event $A_j$ is that clause $C_j$ is not
satisfied, so $\Pr[A_j] = 1/4$. Each clause shares variables with at most
$2(l_j - 1)$ other clauses, so the dependency degree is at most $d = 2m$.
The condition $ep(d+1) \leq 1$ requires $m \leq \lfloor e^{-1}/(2p) \rfloor$,
which limits the density. For sparse formulas, the Local Lemma guarantees
satisfiability.
\end{example}

\begin{example}[Hypergraph Coloring]
\label{ex:hypergraph}
Consider a 3-uniform hypergraph $H = (V, E)$ with $|V| = n$ and
$|E| = m$ edges. We wish to 2-color the vertices so that no edge is
monochromatic. Color each vertex independently red or blue with probability
$1/2$. For each edge $e \in E$, let $A_e$ be the event that $e$ is
monochromatic. Then $\Pr[A_e] = 2 \cdot 2^{-3} = 1/4$. Each edge shares
vertices with at most $3(n-3) = 3n - 9$ other edges (since each of its 3
vertices can belong to at most $n - 3$ other edges). The Local Lemma
condition $ep(d+1) \leq 1$ becomes $e(n/4)(3n - 8) \leq 1$, which is
satisfied for $n \leq O(1)$. More generally, for $k$-uniform hypergraphs,
the condition $ep(d+1) \leq 1$ gives bounds of the form $m \leq
2^{k-1} / (ek)$, ensuring a proper 2-coloring exists.
\end{example}

\begin{example}[$k$-SAT]
For $k$-SAT formulas with clause density $m/n \leq 2^{k-1}/(ek)$,
the Local Lemma guarantees that a random assignment satisfies all clauses
with positive probability. This gives a satisfiability threshold that
matches the upper bound from the DPLL algorithm for small $k$.
\end{example}

%======================================================================
\subsection{Algorithmic Lov\'asz Local Lemma}
\label{sec:algorithmic-lll}
%======================================================================

The Lov\'asz Local Lemma as presented above is a pure existence
result. It tells us that some outcome avoids all bad events but gives
no method to find it. The algorithmic LLL of Moser and
Tardos~\cite{mosertardos2010} addresses precisely this: it provides a
simple, practical algorithm that almost surely finds such an outcome.

\subsubsection*{The variable framework}
Moser--Tardos assumes a probability space defined by a set of
independent random variables $\mathcal{X} = \{X_1,\ldots,X_n\}$. Each
bad event $A_i$ depends on a subset $\text{vbl}(A_i) \subseteq
\mathcal{X}$; it is determined solely by the values of those variables.

The algorithm itself is remarkably simple:

\begin{algorithm}[H]
\caption{Moser--Tardos Resampling}
\label{alg:mosertardos}
\begin{algorithmic}[1]
\State Sample all variables independently from their distributions.
\Loop
   \If{no bad event holds}
      \State \textbf{break}
   \Else
      \State Pick any bad event $A_i$
      \ForAll{$X \in \textnormal{vbl}(A_i)$}
         \State Re-sample $X$ from its distribution
      \EndFor
   \EndIf
\EndLoop
\end{algorithmic}
\end{algorithm}

We start with a random outcome. Whenever a bad event happens, we
resample all the variables it depends on. Because the new values are
fresh independent draws, the process has no memory of the previous
assignment of those variables.

\begin{theorem}[Moser--Tardos, 2010]
\label{thm:algorithmic-lll}
If the events $A_1,\ldots,A_m$ satisfy the LLL condition
$e\,p\,(d+1) \le 1$ (symmetric form), then the Resampling algorithm
terminates with probability~1 at a state where no bad event holds. The
expected number of resampling steps is finite, and is bounded by
$O\!\left(\sum_i \frac{x_i}{1-x_i}\right)$ in the asymmetric
formulation.
\end{theorem}

\begin{proof}[Proof sketch]
We describe the key combinatorial idea. Consider the sequence of events
that are resampled. Associate with each resampled event $A_i$ a
\emph{resampling tree} whose nodes are labeled with the identity of the
bad event and whose edges connect events that are resampled within a
dependency step. Because each resampling chooses fresh values for the
variables of $A_i$, and any event $A_j$ that is not a neighbor of $A_i$
in the dependency graph is unaffected, this tree is a Galton-Watson
process. The LLL condition ensures that the branching factor of this
process is less than~1, so the total size of the tree (and thus the
total number of resamplings) is finite with probability~$1$. The
formalization of this argument bounds the expected resampling count as
$\sum_i \frac{x_i}{1-x_i}$ and requires a careful probabilistic
estimate that we omit here.
\end{proof}

\begin{corollary}[LLL is now constructive]
\label{cor:Moser-Tardos-constructive}
Any existence result that can be established via the Local Lemma under
the standard LLL condition yields an efficient algorithm to find the
desired object whenever the LLL condition holds.
\end{corollary}

\subsubsection*{Examples}
The Algorithmic LLL applies to a far wider range of examples than the
few we list here.

\begin{itemize}
\item {\bf Markovian SAT:} consider a $k$-SAT formula $F$ on $m$ clauses
with at most $d$ dependencies per clause, $e p (d+1) \le 1$.
The Resampling algorithm runs in expected polynomial time.

\item {\bf Packing integer programs:} Suppose we have $m$ objects each with a
penalty probability $p$ and each object depends on $d$ others. Whenever all
objects satisfy the LLL condition, the algorithm finds a solution.

\item {\bf Hypergraph coloring:} For a $k$-uniform hypergraph with edges, the algorithm will find a 2-coloring.

\end{itemize}

\begin{cpplisting}
\label{code:algorithmic-lll}
\begin{lstlisting}[language=C++, frame=single, framerule=0.4pt,
  rulecolor=\color{gray}, backgroundcolor=\color{gray!5},
  basicstyle=\ttfamily\small,
  keywordstyle=\bfseries\color{black},
  commentstyle=\color{gray}, numbers=left,
  numberstyle=\tiny\color{gray}, numbersep=8pt,
  tabsize=2, showstringspaces=false]
#include <vector>
#include <random>
#include <iostream>

// Returns a resampled integer in [0, max_val) uniformly.
int sample_unif(std::mt19937& rng, int max_val) {
  std::uniform_int_distribution<int> dist(0, max_val-1);
  return dist(rng);
}
// Structure of a bad event.
struct Event {
  int id;
  std::vector<int> dependencies; // indices of involved variables
};

// Check if a particular event is currently true.
bool is_event_true(int event_id, const std::vector<int>& assignment,
                   const std::vector<Event>& events) {
  const auto& e = events[event_id];
  // Application-specific logic goes here.
  // In the hypergraph case, check whether the three vertices
  // share the same color.
  return false;
}
// Core resampling function.
std::vector<int> find_good_assignment(
    int n_vars, int n_events,
    const std::vector<Event>& events,
    std::mt19937& rng,
    int max_val) {
  // Initialize randomly
  std::vector<int> assignment(n_vars);
  for (int i = 0; i < n_vars; ++i)
    assignment[i] = sample_unif(rng, max_val);
  // Resampling loop
  bool done = false;
  while (!done) {
    done = true;
    for (int i = 0; i < n_events; ++i) {
      if (is_event_true(i, assignment, events)) {
        // Resample dependent variables
        for (int v : events[i].dependencies) {
          assignment[v] = sample_unif(rng, max_val);
        }
        done = false;
      }
    }
  }
  return assignment;
}
\end{lstlisting}
\end{cpplisting}

\subsubsection*{Beyond the Basic LLL}
The resampling approach generalizes to the lopsided Local Lemma and to
the so-called general LLL, but these are beyond our scope.

%======================================================================
\section{The Method of Conditional Probabilities}
\label{sec:conditional}
%======================================================================

The probabilistic method proves existence but does not construct the
object. The \emph{method of conditional probabilities} transforms a
probabilistic existence argument into a deterministic algorithm by
replacing random choices with deterministic ones that maintain
non-negative conditional expectation.

The key insight is as follows. Suppose we want to show that a random
variable $X$ satisfies $\mathbb{E}[X] \geq 0$, which implies the existence
of an outcome with $X \geq 0$. We can derandomize by setting the variables
$x_1, x_2, \ldots, x_n$ one at a time. At step $i$, we compute the
conditional expectation $\mathbb{E}[X \mid x_1, \ldots, x_{i-1}]$ as a
function of $x_i$, and choose the value of $x_i$ that maximizes this
conditional expectation. The tower property of conditional expectation
ensures that the final value satisfies $X \geq \mathbb{E}[X] \geq 0$.

\begin{theorem}[Derandomization of MAX-3SAT]
\label{thm:derand-maxsat}
For any 3-CNF formula with $m$ clauses, there exists a deterministic
polynomial-time algorithm that finds an assignment satisfying at least
$7m/8$ clauses.
\end{theorem}

\begin{proof}
We apply the method of conditional probabilities to the random variable
$X = \sum_{j=1}^{m} X_j$, where $X_j$ is the indicator for clause $C_j$
being satisfied. The expected value is $\mathbb{E}[X] = 7m/8$.

Initialize all variables as unset. For $i = 1, 2, \ldots, n$, we set
$x_i$ to the value ($0$ or $1$) that maximizes the conditional expectation
$\mathbb{E}[X \mid x_1, \ldots, x_i]$. This conditional expectation can
be computed in polynomial time: for each clause $C_j$, if all its literals
are already determined and none is true, then $X_j = 0$; if any literal
is true, then $X_j = 1$; otherwise, $X_j$ is a Bernoulli random variable
with $\mathbb{E}[X_j \mid \text{partial assignment}] = 1 - 2^{-r}$, where
$r$ is the number of unset literals in $C_j$.

By the tower property,
\[
\mathbb{E}[X \mid x_1, \ldots, x_n] = X(x_1, \ldots, x_n) \geq
\mathbb{E}[X \mid x_1] \geq \cdots \geq \mathbb{E}[X] = \frac{7m}{8}.
\]
Thus the final assignment satisfies at least $7m/8$ clauses.
\end{proof}

The method of conditional probabilities is the ``algorithmic'' backbone of
the probabilistic method. It transforms existential proofs into efficient
algorithms, demonstrating that the power of randomization can often be
eliminated without sacrificing the quality of the result.

\begin{cpplisting}
\begin{lstlisting}[caption={Deterministic MAX-SAT via conditional expectations (C++).}]
#include <vector>
#include <cmath>
#include <algorithm>
using std::vector;

struct MaxSATResult {
  vector<int> assignment;
  int satisfied;
};

// Compute conditional expectation of #satisfied clauses
// given a partial assignment (-1 = unset).
double conditional_expectation(
    const vector<vector<int>>& clauses,
    const vector<int>& asgn) {
  double total = 0.0;
  for (const auto& clause : clauses) {
    bool sat = false;
    int unset = 0;
    for (int lit : clause) {
      int var = std::abs(lit) - 1;       // 0-indexed
      if (asgn[var] == -1) { ++unset; }
      else if ((lit > 0 && asgn[var] == 1) ||
               (lit < 0 && asgn[var] == 0))
        { sat = true; break; }
    }
    if (sat) total += 1.0;
    else if (unset > 0)
      total += 1.0 - std::pow(0.5, unset);
  }
  return total;
}

MaxSATResult deterministic_maxsat(
    const vector<vector<int>>& clauses, int n_vars) {
  vector<int> asgn(n_vars, -1);

  for (int i = 0; i < n_vars; ++i) {
    // Try x_i = 0
    asgn[i] = 0;
    double e0 = conditional_expectation(clauses, asgn);
    // Try x_i = 1
    asgn[i] = 1;
    double e1 = conditional_expectation(clauses, asgn);
    // Choose the value that maximizes E[#satisfied]
    asgn[i] = (e1 >= e0) ? 1 : 0;
  }

  // Count actual satisfied clauses
  int sat = 0;
  for (const auto& clause : clauses) {
    for (int lit : clause) {
      int var = std::abs(lit) - 1;
      if ((lit > 0 && asgn[var] == 1) ||
          (lit < 0 && asgn[var] == 0))
        { ++sat; break; }
    }
  }
  return {asgn, sat};
}
\end{lstlisting}
\end{cpplisting}

%======================================================================
\addcontentsline{toc}{section}{Notes}
\section*{Notes}
%======================================================================

The probabilistic method was introduced by Erd\H{o}s~\cite{erdos1947} in
his 1947 proof of a lower bound on Ramsey numbers. The elegant simplicity
of the approach---``if the average is good, something good must
exist''---has made it one of the most widely used tools in combinatorics.

The Lov\'{a}sz Local Lemma was proved by Erd\H{o}s and
Lov\'{a}sz~\cite{lovasz1975} in 1975. Spencer's 1975
survey~\cite{spencer1975} and the definitive monograph by Alon and
Spencer~\cite{alonspencer1992} (The Probabilistic Method, Wiley) are
standard references.

The method of conditional probabilities was developed by Spencer and
others as a systematic derandomization technique. The algorithmic LLL of
Moser and Tardos~\cite{mosertardos2010} provides an efficient constructive
version of the Local Lemma, resolving a long-standing algorithmic question.

Expander graphs have deep connections to coding theory, pseudorandomness,
and computational complexity. The existence of Ramanujan graphs was
established by Lubotzky--Phillips--Sarnak and Margulis using algebraic
methods; random regular graphs achieve near-optimal expansion with high
probability.

%======================================================================
\addcontentsline{toc}{section}{Problems}
\section*{Problems}
%======================================================================

\begin{problem}
Let $G = G(n, 1/2)$ be the Erd\H{o}s--R\'{e}nyi random graph. Using the
probabilistic method, prove that $G$ contains a triangle-free induced
subgraph on at least $n/4$ vertices.
\end{problem}

\begin{problem}
Prove that for any graph $G$ on $n$ vertices with average degree $\bar{d}$,
there exists an independent set of size at least $\frac{n}{\bar{d} + 1}$.
Use the probabilistic method with a non-uniform distribution.
\end{problem}

\begin{problem}
Let $\phi$ be a CNF formula where each clause has at least $k$ literals.
Show that a random assignment satisfies at least $\prod_{j=1}^{m}
(1 - 2^{-l_j})$ clauses in expectation, where $l_j$ is the length of
clause $j$. Derive a lower bound when $k = 4$.
\end{problem}

\begin{problem}
Use the Lov\'{a}sz Local Lemma to prove that the edges of $K_n$ can be
3-colored so that no monochromatic triangle exists, provided $n \leq
\lfloor e^{-1} \cdot 2^2 / 3 \rfloor$. Compare with the bound obtained
from the union bound.
\end{problem}

\begin{problem}
Consider a $k$-uniform hypergraph $H$ with $m$ edges and maximum degree
$\Delta$ (the maximum number of edges containing any single vertex). Show
that $H$ admits a 2-coloring with no monochromatic edge if $m \leq
2^{k-1} \Delta / k$.
\end{problem}

\begin{problem}
Apply the method of conditional probabilities to the Ramsey problem:
deterministically find a 2-coloring of $K_5$ that minimizes the number
of monochromatic triangles. Write pseudocode for the algorithm.
\end{problem}

\begin{problem}
Let $A_1, \ldots, A_m$ be events each with probability at most $p$, and
suppose each event depends on at most $d$ others. Show that the symmetric
Local Lemma condition $ep(d+1) \leq 1$ implies
$\Pr[\overline{A_1} \cap \cdots \cap \overline{A_m}] \geq
\left(1 - \frac{1}{d+1}\right)^m$.
\end{problem}

\begin{problem}
Construct a random $3$-regular graph on $n = 100$ vertices and compute
its second-largest eigenvalue $|\lambda_2|$ using the power method.
Compare with the Ramanujan bound $2\sqrt{2} \approx 2.828$. Repeat for
100 trials and report the distribution of $|\lambda_2|$.
\end{problem}

\begin{problem}
Show that for a 3-CNF formula with $m$ clauses, the method of conditional
probabilities can be implemented in $O(mn)$ time, where $n$ is the number
of variables.
\end{problem}

\begin{problem}
\textbf{(Erd\H{o}s--Selfridge.)} Let $\mathcal{F}$ be a family of subsets
of $\{1, \ldots, n\}$, each of size at least $k$. If $\sum_{S \in
\mathcal{F}} 2^{-|S|} < 1/2$, then there exists a 2-coloring of
$\{1, \ldots, n\}$ with no monochromatic set in $\mathcal{F}$. Prove this
using the asymmetric Local Lemma, and compare the bound with the union
bound.
\end{problem}

\end{document}
